To systematically address the extension to higher spins, we adopt the Bargmann-Wigner (BW) formalism, which provides a rigorous group-theoretical framework for constructing relativistic wave equations for particles of arbitrary spin $s$. Established by V. Bargmann and E. P. Wigner \cite{bargmann1997a}, this method posits that the wavefunction of a particle with spin $s$ can be represented by a multispinor field $\Psi_{\alpha_1 \alpha_2 \dots \alpha_{2s}}$ of rank $2s$, which effectively describes a composite state of $2s$ fundamental spin-1/2 fermions. To ensure that this object transforms irreducibly under the Poincar\'{e} group, the multispinor must satisfy two conditions: it must be totally symmetric under the exchange of any pair of spinor indices, and it must satisfy the Dirac equation for each index independently:

\begin{equation}
    \left( i\,\gamma^\mu \partial_\mu - m \right)_{\alpha_k \beta}\,
    \Psi_{\alpha_1 \dots \beta \dots \alpha_{2s}}(x)
    = 0,
    \qquad \text{for } k = 1, \dots, 2s .
\end{equation}
To address the extension to higher spins, we adopt the methodology established by Watson in their generalization of the Proca equation \cite{Watson2022}. Their approach serves as the guiding framework for this analysis. Specifically, for the vector boson case ($s=1$), the formalism begins with the Bargmann-Wigner equations for a symmetric rank-2 spinor $\Psi_{\alpha\beta}$. Crucially, the standard scalar mass $m$ is replaced by the chiral mass operator $\hat{M}$ derived from the fermionic theory, applied to each spinor index independently:
\begin{equation}
    (i \gamma^\mu \partial_\mu - \hat{M})_{\alpha \lambda} \Psi_{\lambda \beta}(x) = 0, \quad \text{where} \quad \hat{M} = m e^{-i\alpha\gamma^5}.
\end{equation}
The wavefunction is then expanded in terms of the Clifford algebra generators. However, the symmetry constraint intrinsic to the Bargmann-Wigner formalism ($\Psi_{\alpha\beta} = \Psi_{\beta\alpha}$) strictly eliminates the scalar, pseudoscalar, and axial-vector components. Consequently, the wavefunction reduces naturally to a superposition of only the vector ($A_\mu$) and tensor ($F_{\mu\nu}$) terms:
\begin{equation}
    \Psi(x) = \left( \gamma^\mu A_\mu(x) + \frac{1}{2}\sigma^{\mu\nu} F_{\mu\nu}(x) \right) C.
\end{equation}
Substituting this ansatz back into the fundamental chiral equations leads to a decoupling of the components. While the standard theory yields the Proca equation $\partial_\mu F^{\mu\nu} + m^2 A^\nu = 0$, the chiral extension results in the Generalized Proca Equation:
\begin{equation}
    \partial_\mu F^{\mu\nu} + m^2 \cos^2(\alpha) A^\nu + \dots = 0.
\end{equation}

It is important to clarify that the appearance of the $\cos^2(\alpha)$ factor does not violate the relativistic structure of the theory; rather, it acts as a scaling parameter that defines an effective physical mass $M_{\text{eff}} = m \cos(\alpha)$. This implies that the chiral angle can modulate the observable mass of the boson without altering the underlying geometric consistency of the Poincaré group. Building on this success, we propose to use this established procedure to rigorously examine the pertinency of these chiral degrees of freedom in higher-spin sectors, such as the Rarita-Schwinger field ($s=3/2$).


\section*{Extension to the Rarita-Schwinger Field ($s=3/2$)}

Extending the chiral formalism to the Rarita-Schwinger field ($s=3/2$), which describes massive gravitinos or higher-spin resonances, requires adapting the vector-spinor representation $\psi_\mu$. Following the logic established for the Proca field \cite{Watson2022}, we treat the Rarita-Schwinger field as a composite object transforming under the representation $1 \otimes \frac{1}{2}$ of the Lorentz group.

The fundamental dynamical equation is derived by applying the chiral mass operator $\hat{M}$ to the spinor index of the field, while preserving the vector index structure. We propose the Chiral Rarita-Schwinger Equation as follows:

\begin{equation}
    \epsilon^{\mu\nu\rho\sigma} \gamma^5 \gamma_\nu \partial_\rho \psi_\sigma + i \hat{M} \sigma^{\mu\nu} \psi_\nu = 0,
\end{equation}

where $\psi_\mu$ is the vector-spinor field. Consistently with the $s=1$ case, the chiral mass operator $\hat{M}$ acts on the spinor components to project the effective mass. Assuming a symmetric alignment of the chiral bases inherent to the composite representation, the mass term generalizes to:

\begin{equation}
    \hat{M} \psi_\nu \rightarrow m \cos^2(\alpha) \psi_\nu.
\end{equation}

This substitution leads to a modified equation of motion that explicitly includes the chiral scaling factor:

\begin{equation}
    (i \gamma^\mu \partial_\mu - m \cos^2(\alpha)) \psi_\nu - \dots = 0.
\end{equation}

The presence of the $\cos^2(\alpha)$ factor here is physically significant. It suggests that for a spin-$3/2$ particle, the "screening" of the mass by the chiral angle follows the same geometric projection principle observed in the vector boson sector, reinforcing the universality of the effective mass $M_{\text{eff}} = m \cos(\alpha)$ across higher-spin representations.

\section*{Discussion on Generalization via Projection Operators}

The explicit construction of equations for higher spins (like $s=3/2$ or $s=2$) using standard component notation becomes exponentially complex. To solve this, Watson and Musielak propose a powerful generalization based on Orthogonal Idempotents.

Instead of manually constructing the Lagrangians, this method utilizes projection operators $\hat{P}$ to define the physical subspaces directly. For a system of $N$ constituent spinors (which would form a spin-$N/2$ particle), the generalized projection operator is defined as the tensor product of the individual helicity projectors:

\begin{equation}
    \hat{P}_{s_1, s_2, \dots, s_N}(\hat{n}) = \hat{P}_{s_1}(\hat{n}) \otimes \hat{P}_{s_2}(\hat{n}) \otimes \dots \otimes \hat{P}_{s_N}(\hat{n}),
\end{equation}

where each $\hat{P}_{s_i}(\hat{n})$ projects the spin state relative to a quantization axis $\hat{n}$. For the specific case of the Rarita-Schwinger field ($N=3$), one selects the symmetric subspace of three entangled spinors.

This projector formalism is superior because it automatically handles the chiral constraints. By defining the chiral basis for each projector (parameterized by angles $\alpha_1, \alpha_2, \dots$), the resulting equation of motion emerges naturally with the mass terms already modulated by the relative chiral angles. Thus, the effective mass $M_{\text{eff}}$ is not an ad-hoc insertion but an eigenvalue of the generalized projection operator acting on the vacuum state:

\begin{equation}
    \hat{P}_{\text{total}} \Psi = M_{\text{eff}} \Psi.
\end{equation}

This confirms that the chiral degrees of freedom are fundamental to the algebraic structure of the Poincaré group representations, providing a rigorous path to derive wave equations for any arbitrary spin $s$.
